\documentclass[../../../main.tex]{subfiles}
\begin{document}

% Keep this file in sync with the matching problem_07_review.md sidecar.

\problemsetup
  {TODO: Concise Problem Title (contributor\_01 / problem\_07)}
  {TODO: Mathematics / Statistics / Data Science}
  {TODO: subtopic}
  {Qualifying-exam-plus}
  {Zhixian Zhang}
  {TODO: contributor contact}
  {Draft}
\problemtags
  {TODO: compute / prove / diagnose / simulate / plan / debug}
  {TODO: challenged capability}
  {TODO: exact match / numerical tolerance / structural check / rubric-based / agent-as-judge / programmatic test}
  {TODO: short observed failure mode}
\begin{msbproblem}

\begin{problemstatement}
\subsection*{Setting and Definitions} 
Let $\Theta \subseteq \mathbb{R}^d$ be a closed and convex set of model parameters. Deploying a model $f_\theta$ induces a data distribution $\mathcal{D}(\theta)$. For any $w \in \Theta$, the risk of a model $f_w$ is defined as:
\[
\mathcal{L}(\theta; w) = \mathbb{E}_{Z \sim \mathcal{D}(w)} [\ell(Z; \theta)]
\]
where $\ell(z; \cdot)$ is the loss function. 
Consider the sequence generated by the successive applications of the retraining mapping, defined as $\theta_{t+1} = T(\theta_t)$, where the mapping $T: \Theta \rightarrow \Theta$ is given by:
\[
T(w) = \arg\min_{\theta \in \Theta} \mathcal{L}(\theta; w)
\]
We denote the empirical version of the retraining mapping as $\hat{T}(w)$, yielding the sample-based trajectory, starting from a given $\theta_0 = \hat \theta_0$:
\[
\hat{\theta}_t = \widehat{T}(\hat{\theta}_{t-1}) = \arg\min_{w \in \Theta} \frac{1}{n} \sum_{i=1}^n \ell(Z_{i,t}; w), \quad \text{where } Z_{i,t} \sim \mathcal{D}(\hat{\theta}_{t-1}) \text{ i.i.d.}
\]

\subsection*{Assumption}
Assume the following conditions hold for some constants $\gamma, \beta, \epsilon > 0$:
\begin{enumerate}
    \item For any fixed data point $z$, the loss function $\ell(z; \cdot)$ is $\gamma$-strongly convex.
    \item The gradient mapping $\nabla_\theta \ell(z; \theta)$ is $\beta$-Lipschitz continuous in $z$ and $\theta$.
    \item The distribution mapping $\mathcal{D}(\cdot)$ is $\epsilon$-sensitive under the Wasserstein-1 distance.
    \item For each $\theta_t$, there exists a local neighborhood $C$ of $\theta_t$ such that the loss function $\ell(z; \theta)$ is $L_C(z)$-Lipschitz continuous with respect to $\theta$ for all $z$ with $\mathbb{E}_{Z \sim \mathcal{D}(\theta)}[L(Z)]^2 < \infty$.
    \item The conditional variance of the loss gradient is positive definite:
    \[
    V_{w}(\theta) = \text{Var}_{Z \sim \mathcal{D}(w)} \left[ \nabla_\theta \ell(Z; \theta) \right] \succ 0
    \]
    %\item The risk function $\mathcal{L}(\theta; w)$ is differentiable on $\theta$.
    \item For any sample size $n$ and iteration $t$, the estimator $\widehat{\theta}_t$ has a well-defined density function with respect to the Lebesgue measure, and its corresponding characteristic function is absolutely integrable over $\mathbb{R}^d$.
\end{enumerate}

\subsection*{Question 1}
Let $\kappa = \frac{\epsilon\beta}{\gamma}$. Suppose we execute the retrainig to generate the sequence $\{\theta_{t}\}_{t=1}$.
\begin{enumerate}
    \item Prove that for any $\theta, \theta' \in \Theta$:
    \[
    \|T(\theta) - T(\theta')\|_2 \le \frac{\epsilon\beta}{\gamma} \|\theta - \theta'\|_2.
    \]
    \item Does the sequence $\theta_t$ converge under the condition $\kappa < 1$? If convergence is guaranteed, prove that the sequence converges to its limit point $\theta^*$ at a linear rate in the $\ell_2$-norm, and derive the minimum number of iterations $t$ required to ensure that $\|\theta_t - \theta^*\|_2 \le \delta$.
\end{enumerate}

\subsection*{Question 2}
Let $\Theta \subseteq \mathbb{R}^d$ be the closed and convex set of model parameters. Suppose the condition $\kappa < 1$ holds.
\begin{enumerate}
    \item For any fixed iteration $t \ge 1$, does $\widehat{\theta}_t \xrightarrow{P} \theta_t$ hold? Provide a rigorous and detailed justification.
    \item For any fixed iteration $t \ge 1$, does the empirical estimator $\widehat{\theta}_t$ exhibit asymptotic normality to $\theta_t$ as $n \to \infty$? If so, prove this property rigorously in detail and show the asymptotic covariance matrix.
\end{enumerate}


\end{problemstatement}

\begin{answerkey}
\subsection*{Solution to Question 1}

\begin{enumerate}
    \item By the first-order optimality conditions of the retraining mapping, we have:
\[
\nabla_\theta \mathcal{L}(T(\theta); \theta) = 0 \quad \text{and} \quad \nabla_\theta \mathcal{L}(T(\theta'); \theta') = 0.
\]
Applying the $\gamma$-strong convexity of $\mathcal{L}(\cdot; \theta)$ followed by the Kantorovich-Rubinstein theorem for the $\beta$-Lipschitz gradient $f(z) = \nabla_\theta \ell(z; T(\theta'))$, we obtain:
\begin{align*}
\gamma \|T(\theta) - T(\theta')\|_2 & \le \| \nabla_\theta \mathcal{L}(T(\theta); \theta) - \nabla_\theta \mathcal{L}(T(\theta'); \theta) \|_2 \\
& = \| \nabla_\theta \mathcal{L}(T(\theta'); \theta') - \nabla_\theta \mathcal{L}(T(\theta'); \theta) \|_2 \\
& = \left\| \mathbb{E}_{Z \sim \mathcal{D}(\theta')}[\nabla_\theta \ell(Z; T(\theta'))] - \mathbb{E}_{Z \sim \mathcal{D}(\theta)}[\nabla_\theta \ell(Z; T(\theta'))] \right\|_2 \\
& \le \beta \cdot \mathcal{W}_1\big(\mathcal{D}(\theta), \mathcal{D}(\theta')\big).
\end{align*}
Invoking the $\epsilon$-sensitivity condition $\mathcal{W}_1\big(\mathcal{D}(\theta), \mathcal{D}(\theta')\big) \le \epsilon \|\theta - \theta'\|_2$ yields:
\[
\gamma \|T(\theta) - T(\theta')\|_2 \le \epsilon\beta \|\theta - \theta'\|_2 \implies \|T(\theta) - T(\theta')\|_2 \le \frac{\epsilon\beta}{\gamma} \|\theta - \theta'\|_2.
\]

    \item Since $\kappa < 1$ and $\Theta \subseteq \mathbb{R}^d$ is closed, $T(\cdot)$ is a strict contraction. By the \emph{Banach Fixed-Point Theorem}, there exists a unique fixed point $\theta^* = T(\theta^*)$. 

By the contraction property and mathematical induction, we have:
\[
\|\theta_t - \theta^*\|_2 = \|T(\theta_{t-1}) - T(\theta^*)\|_2 \le \kappa \|\theta_{t-1} - \theta^*\|_2 \le \kappa^t \|\theta_0 - \theta^*\|_2.
\]
Setting $\kappa^t \|\theta_0 - \theta^*\|_2 \le \delta$ and solving for $t$ immediately yields the linear convergence iteration complexity:
\[
t \ge \frac{\ln\left(\|\theta_0 - \theta^*\|_2 / \delta\right)}{\ln(1/\kappa)}.
\]
\end{enumerate}


\subsection*{Solution to Question 2}

\begin{enumerate}
    \item We have
\[
\begin{aligned}
    \|\theta_t-\widehat\theta_t\|
    &=
    \|T(\theta_{t-1})-\widehat T(\widehat\theta_{t-1})\| \\
    &\le
    \|T(\widehat\theta_{t-1})-\widehat T(\widehat\theta_{t-1})\|
    +
    \|T(\theta_{t-1})-T(\widehat\theta_{t-1})\| \\
    &\le
    \|T(\widehat\theta_{t-1})-\widehat T(\widehat\theta_{t-1})\|
    +
    \frac{\varepsilon\beta}{\gamma}
    \|\theta_{t-1}-\widehat\theta_{t-1}\|,
\end{aligned}
\]
By the local Lipschitz condition, there exists $\varepsilon_0>0$ such that
\[
    \sup_{\theta:\|\theta-T(\widehat\theta_{t-1})\|\le\varepsilon_0}
    \left|
        \widehat{\mathcal{L}}(\theta,\widehat\theta_{t-1})
        -
        \mathcal{L}(\theta,\widehat\theta_{t-1})
    \right|
    \xrightarrow{P}0.
\]
Since $\ell$ is strongly convex for any $\theta$,
$T(\widehat\theta_{t-1})$ is unique, so there exists $\delta$
such that $\mathcal{L}(\theta,\widehat\theta_{t-1})-\mathcal{L}(T(\widehat\theta_{t-1}),\widehat\theta_{t-1}) \geq \delta$
for all
$\theta\in\{\theta:\|\theta-T(\widehat\theta_{t-1})\|=\varepsilon_0\}$.
Then it follows that
\[
\begin{aligned}
&\inf_{\|\theta-T(\widehat\theta_{t-1})\|=\varepsilon_0}
    \widehat{\mathcal{L}}(\theta,\widehat\theta_{t-1})
    -
    \widehat{\mathcal{L}}(T(\widehat\theta_{t-1}),\widehat\theta_{t-1}) \\
=&
\inf_{\|\theta-T(\widehat\theta_{t-1})\|=\varepsilon_0}
\Big[
    \bigl(
        \widehat{\mathcal{L}}(\theta,\widehat\theta_{t-1})
        -
        \mathcal{L}(\theta,\widehat\theta_{t-1})
    \bigr)
    +
    \bigl(
        \mathcal{L}(\theta,\widehat\theta_{t-1})
        -
        \mathcal{L}(T(\widehat\theta_{t-1}),\widehat\theta_{t-1})
    \bigr) \\
&\hspace{4.5cm}
    +
    \bigl(
        \mathcal{L}(T(\widehat\theta_{t-1}),\widehat\theta_{t-1})
        -
        \widehat{\mathcal{L}}(T(\widehat\theta_{t-1}),\widehat\theta_{t-1})
    \bigr)
\Big] \\
\ge& \delta-o_P(1).
\end{aligned}
\]
We consider any fixed $\theta$ such that
$\|\theta-T(\widehat\theta_{t-1})\|\ge\varepsilon_0$. Let $\omega=\frac{\theta-T(\widehat\theta_{t-1})}{\|\theta-T(\widehat\theta_{t-1})\|}\varepsilon_0+T(\widehat\theta_{t-1})$, by the convexity of $\widehat{\mathcal{L}}$, we have
\[
\begin{aligned}
&
\widehat{\mathcal{L}}(\theta,\widehat\theta_{t-1})
-
\widehat{\mathcal{L}}(T(\widehat\theta_{t-1}),\widehat\theta_{t-1}) \\
\ge&
    \frac{\|\theta-T(\widehat\theta_{t-1})\|}
    {\|\omega-T(\widehat\theta_{t-1})\|}
    \left(
        \widehat{\mathcal{L}}(\omega,\widehat\theta_{t-1})
        -
        \widehat{\mathcal{L}}(T(\widehat\theta_{t-1}),\widehat\theta_{t-1})
    \right) \\
\ge&
    \frac{\|\theta-T(\widehat\theta_{t-1})\|}{\varepsilon_0}
    \bigl(\delta-o_P(1)\bigr)
    \ge
    \delta-o_P(1),
\end{aligned}
\]
Thus, no $\theta$ such that
$\|\theta-T(\widehat\theta_{t-1})\|\ge\varepsilon_0$ can be the minimizer of
$\widehat{\mathcal{L}}(\theta,\widehat\theta_{t-1})$, so
\[
    \|T(\widehat\theta_{t-1})-\widehat T(\widehat\theta_{t-1})\|
    \xrightarrow{P}0.
\]
For all $t \geq 1$, we have
\[
    \|\widehat\theta_t-\theta_t\|
    \le
    \sum_{i=0}^{t}
    \left(\frac{\varepsilon\beta}{\gamma}\right)^{t-i}
    \|T(\widehat\theta_i)-\widehat T(\widehat\theta_i)\|
    \xrightarrow{P}0.
\]

    \item 
\begin{enumerate}
    \item \textbf{Decomposition.}

    Define the normalized error
    \[
        U_t := \sqrt{n}(\widehat\theta_t-\theta_t),
        \qquad
        \widetilde\theta_t := T(\widehat\theta_{t-1}).
    \]
    The key decomposition is
    \[
        \widehat\theta_t-\theta_t
        =
        \bigl(\widetilde\theta_t-\theta_t\bigr)
        +
        \bigl(\widehat\theta_t-\widetilde\theta_t\bigr).
    \]

    \item \textbf{Conditional distribution.}

    Conditional on $\widehat\theta_{t-1}$, the samples at time $t$ are i.i.d.
    from $D(\widehat\theta_{t-1})$, and
    $\widetilde\theta_t=T(\widehat\theta_{t-1})$ is fixed. For the second term, the proof uses the classical M-estimator CLT, obtained
    by empirical process approximation and a second-order Taylor expansion of
    the empirical risk:
    \[
        \sqrt{n}(\widehat\theta_t-\widetilde\theta_t)
        \mid \widehat\theta_{t-1}
        \xrightarrow{D}
        N\bigl(0,\Sigma_{\widehat\theta_{t-1}}(\widetilde\theta_t)\bigr).
    \]

    Since
    $\widehat\theta_{t-1}=\theta_{t-1}+U_{t-1}/\sqrt{n}$, this gives the
    conditional limiting distribution
    \[
        U_t\mid U_{t-1}
        \xrightarrow{D}
        N\bigl(\mu(U_{t-1}),\Sigma(U_{t-1})\bigr),
    \]
    where
    \[
        \mu(U_{t-1})
        =
        \sqrt{n}\left[
        T\left(\theta_{t-1}+\frac{U_{t-1}}{\sqrt{n}}\right)
        -
        T(\theta_{t-1})
        \right].
    \]

    \item \textbf{Marginal distribution.}

    The proof then turns the conditional CLT into a marginal CLT by induction
    and the characteristic function method. By the tower
    property:
    \[
        \phi_{U_t}(s)
        =
        \mathbb{E}\left[
            \phi_{U_t\mid U_{t-1}}(s\mid U_{t-1})
        \right].
    \]
    Therefore, once we know the conditional limiting distribution of
    $U_t\mid U_{t-1}$, we can obtain the marginal characteristic function by averaging over the law of $U_{t-1}$.
    By calculation, the limiting marginal characteristic function is
    Gaussian:
    \[
        U_t\xrightarrow{D}N(0,V_t),
    \]
    with recursive variance
    \[
        V_t
        =
        \nabla T(\theta_{t-1})V_{t-1}\nabla T(\theta_{t-1})^\top
        +
        \Sigma_{\theta_{t-1}}(\theta_t).
    \]
    Here
    \[
        \Sigma_{\theta_{t-1}}(\theta_t)
        =
        H_{\theta_{t-1}}(\theta_t)^{-1}
        V_{\theta_{t-1}}(\theta_t)
        H_{\theta_{t-1}}(\theta_t)^{-1},
    \]
    where
    \begin{align*}
        H_{\theta_{t-1}}(\theta_t)
        &:=
        \mathbb{E}_{z\sim D(\theta_{t-1})} \nabla_\theta^2\ell(z;\theta_t), \\
        V_{\theta_{t-1}}(\theta_t)
        &:=
        \operatorname{Cov}_{z\sim D(\theta_{t-1})}
        \bigl(\nabla_\theta \ell(z;\theta_t)\bigr).
    \end{align*}
    
\end{enumerate}

\end{enumerate}


\end{answerkey}

\begin{gradingscheme}

\begin{enumerate}[label=\arabic*.]
\item \textbf{Question 1(1): Contraction Property and Sensitivity Bounds}
\begin{itemize}
\item \textbf{Criterion 1.1 (Optimality \& Convexity):} The response must correctly invoke the first-order optimality conditions ($\nabla_\theta \mathcal{L} = 0$) and apply the $\gamma$-strong convexity of the loss function to set up the foundational upper bound: $\gamma \|T(\theta) - T(\theta')\|_2 \le \| \nabla_\theta \mathcal{L}(T(\theta); \theta) - \nabla_\theta \mathcal{L}(T(\theta'); \theta) \|_2$.
\item \textbf{Criterion 1.2 (Kantorovich-Rubinstein Alignment):} The response must correctly rewrite the gradient difference as an expectation discrepancy under shifting distributions ($\mathcal{D}(\theta)$ vs. $\mathcal{D}(\theta')$) and apply the Kantorovich-Rubinstein theorem using the $\beta$-Lipschitz property of the gradient to bound it by $\beta \cdot \mathcal{W}_1\big(\mathcal{D}(\theta), \mathcal{D}(\theta')\big)$.
\item \textbf{Criterion 1.3 (Final Contraction Parameter):} The response must inject the $\epsilon$-sensitivity condition to conclude the strict contractive modulus: $\|T(\theta) - T(\theta')\|_2 \le \frac{\epsilon\beta}{\gamma} \|\theta - \theta'\|_2$.
\end{itemize}

\item \textbf{Question 1(2): Fixed-Point Convergence and Iteration Complexity}
\begin{itemize}
\item \textbf{Criterion 2.1 (Banach Application):} The response must explicitly invoke the Banach Fixed-Point Theorem, validating that since $\kappa = \frac{\epsilon\beta}{\gamma} < 1$ and $\Theta$ is a closed subset of $\mathbb{R}^d$, a unique fixed point $\theta^*$ exists.
\item \textbf{Criterion 2.2 (Complexity Derivation):} The response must utilize mathematical induction to bound the tracking error at step $t$ by $\kappa^t \|\theta_0 - \theta^*\|_2$, and algebraically invert the exponential inequality to explicitly derive the logarithmic iteration complexity threshold: $t \ge \frac{\ln\left(\|\theta_0 - \theta^*\|_2 / \delta\right)}{\ln(1/\kappa)}$.
\end{itemize}

\item \textbf{Question 2(1): Probability Consistency under Empirical Risk Perturbations}
\begin{itemize}
\item \textbf{Criterion 3.1 (Triangle Inequality Expansion):} The response must execute a proper multi-step triangle inequality breakdown to isolate the empirical optimization error from the theoretical retraining mapping: $\|\theta_t-\widehat\theta_t\| \le \|T(\widehat\theta_{t-1})-\widehat T(\widehat\theta_{t-1})\| + \kappa\|\theta_{t-1}-\widehat\theta_{t-1}\|$.
\item \textbf{Criterion 3.2 (Convex Compact Boundary Auditing):} The response must demonstrate understanding of M-estimator consistency via convexity. It must construct a compact boundary $\|\theta-T(\widehat\theta_{t-1})\|=\varepsilon_0$, utilize the local Uniform Law of Large Numbers ($\sup |\widehat{\mathcal{L}} - \mathcal{L}| \xrightarrow{P}0$) alongside strong convexity to establish an energy barrier $\delta - o_P(1)$, and use a convex combination projection ($\omega$) to rule out external minimizers.
\item \textbf{Criterion 3.3 (Asymptotic Convergence Mapping):} The response must rigorously show that the point-wise empirical mapping error vanishes in probability ($\|T(\widehat\theta_{t-1})-\widehat T(\widehat\theta_{t-1})\| \xrightarrow{P}0$) and use a finite sum geometric expansion over $i=0$ to $t$ to conclude that the total accumulated error converges to zero in probability ($\|\widehat\theta_t-\theta_t\| \xrightarrow{P}0$).
\end{itemize}

\item \textbf{Question 2(2): Advanced Asymptotic Normality and Tower Property Aggregation}
\begin{itemize}
\item \textbf{Criterion 4.1 (Conditional M-Estimator CLT):} The response must formalize the temporal dependency structure by conditioning on the filtration or past parameters ($\mid \widehat\theta_{t-1}$). It must correctly derive the conditional limiting distribution of the scaled error vector $U_t \mid U_{t-1}$ as a Gaussian distribution driven by the localized drift vector $\mu(U_{t-1})$ and regularized covariance matrix.
\item \textbf{Criterion 4.2 (Characteristic Function Induction):} The response must explicitly deploy the Tower Property of expectation on the joint probability space to transition from a conditional CLT to a marginal CLT, linking the characteristic functions via $\phi_{U_t}(s) = \mathbb{E}\left[\phi_{U_t\mid U_{t-1}}(s\mid U_{t-1})\right]$.
\item \textbf{Criterion 4.3 (Recursive Variance \& Sandwich Formulation):} The response must compute the definitive marginal Gaussian limit $U_t \xrightarrow{D} N(0, V_t)$ and specify the exact discrete-time Riccati-like recursive variance equation: $V_t = \nabla T(\theta_{t-1})V_{t-1}\nabla T(\theta_{t-1})^\top + \Sigma_{\theta_{t-1}}(\theta_t)$, where $\Sigma$ must be formulated as the standard parametric "sandwich" covariance matrix involving the Hessian $H$ and the gradient covariance $V$.
\end{itemize}
\end{enumerate}
\end{gradingscheme}

\begin{editornotes}
Suggested checks:
\begin{itemize}
\item Is every symbol defined?
\item Does the rubric use exact answer / tolerance for objective questions or a structural / judge rubric for proof-style questions?
\item Can the verifier grade deterministically, or at least be audited line by line from the written schema?
\item Does the problem require non-routine reasoning?
\item Does the matching review sidecar contain five raw model responses and a failure-mode summary?
\end{itemize}
\end{editornotes}

\end{msbproblem}

\end{document}
